\documentclass{article}
\usepackage[T1]{fontenc}
\usepackage{lmodern}
\usepackage{graphicx}
\usepackage{amsmath,amssymb}
\usepackage[a4paper,margin=2cm]{geometry}
\usepackage[absolute,overlay]{textpos}
\setlength{\TPHorizModule}{1mm}
\setlength{\TPVertModule}{1mm}
\usepackage[indonesian]{babel}
\usepackage{hyperref}
\setlength{\emergencystretch}{2em}
% unit: Halaman 15

\begin{document}

% === Halaman 15 (llm) ===

"Cryptography lies at the heart of all cybersecuriry technologies, and understanding the role it plays is vital to understanding how to secure cyberspace" (Keith M. Martin)

\textcolor{red}{(Kriptografi merupakan inti dari semua teknologi keamanan siber, dan memahami peran yang dimainkannya sangat penting untuk memahami cara mengamankan ruang siber)}

\vspace{10mm}

Kriptografer masih menjadi sosok yang langka di Indonesia. Namun, bidang ini sangat penting untuk segera ditumbuhkan

\end{document}
